\section{The higher compact current retract}
\label{sec:current-higher-retract}

The compact projection $\Pi=\pi N_Q$ is the linear term of a
symmetric higher morphism.  Its homotopy $H=N_Q^{-1}h_0N_Q$ also
extends to higher arities.  A polynomial path which scales the
contractible variables constructs all the required coefficients.
The points, cutoffs, and compact supports remain fixed along this path.

Use the finite supported current complex $Q$ of
Theorem~\ref{thm:bf-compact-quantum-contraction}, with
$\delta_i=x_i-y_i$, odd variables $\theta_i$, and
$A=\C[y_1,\ldots,y_N][\hbar]$.
Its differential is $D_Q=N_Q^{-1}d_0N_Q$, where
\[
 d_0=\sum_i\delta_i\partial_{\theta_i},\qquad
 N_Q=e^{-\hbar K_Q},\qquad
 K_Q=\sum_{i<j}p_{ij}\partial_{\delta_i}\partial_{\delta_j}.
\]
The symmetric degree-one brackets on $Q$ are $D_Q$ and its
second-order Leibniz defect.  All higher brackets vanish.
Give $A$ the zero differential and zero brackets.
For a degree-zero linear map $T$ between graded commutative algebras,
write $C_n(T)$ for its partition formula
\eqref{eq:current-cumulant-taylor}, with $N$ replaced by $T$.
Thus $C_1(T)=T$ and $C_2(T)(a,b)=T(ab)-T(a)T(b)$.

\begin{proposition}[Higher moment projection]
\label{prop:current-higher-moment-projection}
For the finite list of current jets, the maps
\[
 \Lambda_n=C_n(L_U):\mathcal B_\kappa(U)^{\otimes n}\longrightarrow A
\]
define a graded symmetric $L_\infty[1]$ morphism.
Here $L_U=RN_U$ is the full smooth moment chain map of
Proposition~\ref{prop:bf-compact-full-boundary-split}.
The annular inclusion $\iota_U:A\to\mathcal B_\kappa(U)$ is a
strict morphism and $\Lambda\iota_U=1$ as coalgebra maps.
On the embedded subcomplex $Q$, the Taylor maps are
$\Lambda_n|_Q=C_n(\Pi)$.
This restriction is an $L_\infty[1]$ quasi-isomorphism.
\end{proposition}
\begin{proof}
The moment map $R$ is a classical algebra map and chain map to $A$.
It therefore induces a strict morphism from the abelian complex
$(B,d)$ to the abelian complex $A$.
Compose it with the symmetric Taylor morphism of
Theorem~\ref{thm:current-supported-taylor}.
Since $R$ preserves every product, its composite has coefficients
$R C_n(N_U)=C_n(RN_U)$.
This proves the asserted morphism on the full smooth carrier.

The differential and its binary defect vanish on annular polynomials.
The inclusion is therefore strict.  The identities
$N_U\iota_U=\iota_U$ and $R\iota_U=1$ make its unary composite
the identity.  Partition inversion makes every higher composite zero.
Thus the composite is the identity coalgebra map, including all arities.
On $Q$, the equality $L_U\rho=\Pi$ gives the stated restriction.
Its linear coefficient is the previously proved quasi-isomorphism
$\Pi:Q\to A$.
\end{proof}

The full boundary complex need not retract by a chain homotopy onto
this finite polynomial algebra.  The explicit homotopy below is on
the actual subcomplex $Q$, which contains all the chosen compact
current representatives and their odd primitives.

Let $\Omega=\C[\tau,d\tau]$ be polynomial forms on $[0,1]$, with
$|\tau|=0$, $|d\tau|=1$, and $(d\tau)^2=0$.
The parameter $\tau$ is a homotopy parameter, separate from the
physical interval coordinate and the curve coordinate.
Place the form factor first.  On $\Omega\otimes Q$, use
\[
 (\alpha\otimes q)(\beta\otimes r)
       =(-1)^{|q||\beta|}\alpha\beta\otimes qr,
 \qquad
 D_{\mathrm{tot}}(\alpha\otimes q)
       =d_\tau\alpha\otimes q+(-1)^{|\alpha|}\alpha\otimes D_Qq.
\]
The classical total differential replaces $D_Q$ by $d_0$.
Extend $N_Q$ to $N_{\mathrm{out}}=1\otimes N_Q$.
Define the classical algebra map
\[
 P:Q\longrightarrow\Omega\otimes Q,\qquad
 P(y_i)=y_i,\quad P(\theta_i)=\tau\theta_i,
 \quad P(\delta_i)=\tau\delta_i+d\tau\theta_i.
\]
All displayed expressions have the degree of their input.
Set
\[
 T=N_{\mathrm{out}}^{-1}PN_Q,
 \qquad \mathcal T_n=C_n(T).
\]
Let $j_A:A\to\Omega\otimes A$ include the constant forms,
and let $\Lambda_\Omega=C(1\otimes\Pi)$ be the higher moment
map on the target complex.

\begin{theorem}[Polynomial higher homotopy]
\label{thm:current-polynomial-higher-homotopy}
The maps $\mathcal T_n$ define an $L_\infty[1]$ morphism from
$Q$ to $\Omega\otimes Q$, both with the brackets determined by
their actual quantum differentials.
Endpoint evaluation gives
\[
 \operatorname{ev}_0\mathcal T=\iota\Lambda|_Q,
 \qquad \operatorname{ev}_1\mathcal T=1_Q.
\]
The homotopy fixes the inclusion of $A$ and satisfies
$\Lambda_\Omega\mathcal T=j_A\Lambda|_Q$.
Its unary coefficient recovers the compact chain homotopy:
\[
 \int_0^1 [d\tau]T=H.
\]
Here $[d\tau]$ means the coefficient of the leftmost $d\tau$.
Every formula is polynomial in $\hbar$ on each polynomial input.
After substitution of the compact representatives, all outputs
remain supported in the same disk $U$.
\end{theorem}
\begin{proof}
The classical total differential satisfies
\[
 d_{\mathrm{tot}}P(\theta_i)=d\tau\theta_i+\tau\delta_i
                           =P(d_0\theta_i),\qquad
 d_{\mathrm{tot}}P(\delta_i)=d\tau\delta_i-d\tau\delta_i=0.
\]
It also kills $P(y_i)$.  Both classical differentials are derivations,
so $P$ is a chain map on the whole polynomial algebra.
The Gaussian conjugation commutes with the parameter differential.
Consequently $T$ is a chain map for the quantum differentials.

For any such chain map, the coalgebra construction
$C(T)=E^{-1}S^c(T)E$ intertwines the descendant coderivations.
Apply the partition argument of
Theorem~\ref{thm:current-supported-taylor} to source and target.
The target differential still has order at most two.
It therefore has only the stated unary and binary brackets.
This proves every Taylor identity for $\mathcal T$.

Endpoint evaluation sets $d\tau$ to zero.
It gives $P_0=\iota\pi$ and $P_1=1$.
Since $N_Q^{\pm1}\iota=\iota$, one obtains
$T_0=\iota\Pi$ and $T_1=1$.
Cumulants commute with composition, as follows by cancellation of
the middle $E E^{-1}$ in their coalgebra expressions.
Thus $C(T_0)=\iota C(\Pi)=\iota\Lambda|_Q$,
while $C(T_1)$ is the identity with no higher coefficients.
Also $P\iota=\iota$ and $\pi P=\pi$.
These equations prove the assertions about inclusion and projection
after Gaussian conjugation and the same coalgebra composition.

For a monomial $q$ of positive total degree $r$ in
$\delta_i,\theta_i$, substitution gives
\[
 P(q)=\tau^r q+d\tau\,\tau^{r-1}
                         \sum_i\theta_i\partial_{\delta_i}q.
\]
The left odd factors carry exactly the Koszul sign of the
substitution.  A degree-zero monomial in these variables has no
$d\tau$ coefficient.  Integrating $\tau^{r-1}$ gives $1/r$,
so $\int[d\tau]P=h_0$.  Conjugation then gives
$\int[d\tau]T=N_Q^{-1}h_0N_Q=H$.
No operation changes the chosen compact representatives or their
support.  The Gaussian exponentials terminate, proving the final assertions.
\end{proof}

The first higher coefficients display the correction required during
the homotopy.  The chain rule gives
$K_{\mathrm{out}}P=\tau^2PK_Q$, and therefore
\[
 T=P\exp\bigl(-\hbar(1-\tau^2)K_Q\bigr).
\]
Put $a_i=y_i+\tau\delta_i+d\tau\theta_i=T(x_i)$.
For two distinct current indices,
\[
 \mathcal T_2(x_i,x_j)=-\hbar(1-\tau^2)p_{ij}.
\]
Its total differential is $2\hbar\tau d\tau p_{ij}$.
The target binary bracket on $(a_i,a_j)$ is
$-2\hbar\tau d\tau p_{ij}$.
They cancel in the arity-two morphism equation, whose source
bracket on $(x_i,x_j)$ is zero.
Thus keeping only the linear path would violate that equation at
nonzero covariance and an interior value of $\tau$.

The first connected three-input coefficient is
\[
 \mathcal T_3(x_1^2,x_2,x_3)
       =2\hbar^2(1-\tau^2)^2p_{12}p_{13}.
\]
It equals the moment-projection coefficient at $\tau=0$ and
vanishes at $\tau=1$.  The full polynomial path supplies its
compatibility with every other Taylor coefficient.

The higher moment map extends the same $L_U$ used by
Proposition~\ref{prop:interval-supported-current-chain-map}.
Its restriction supplies a symmetric higher inverse to the annular
current inclusion, and the polynomial path refines the same compact
homotopy used in the supported product comparison.
The ordered interval operations and distributional collision domains
still require their own compatible maps.
